% Generated draft: requires mathematical and source review.
\documentclass[11pt]{article}
\usepackage[margin=25mm]{geometry}
\usepackage{fontspec}
\setmainfont{Latin Modern Roman}
\usepackage{amsmath,amssymb,mathtools,unicode-math}
\setmathfont{Latin Modern Math}
\usepackage{enumitem,needspace,xurl,hyperref}
\setlength{\parindent}{0pt}
\setlength{\parskip}{5pt}
\setlength{\emergencystretch}{3em}
\newcommand{\N}{\mathbb{N}}
\newcommand{\Z}{\mathbb{Z}}
\newcommand{\Q}{\mathbb{Q}}
\newcommand{\R}{\mathbb{R}}
\newcommand{\C}{\mathbb{C}}
\newcommand{\F}{\mathbb{F}}
\newcommand{\A}{\mathbb{A}}
\newcommand{\PP}{\mathbb{P}}
\newcommand{\E}{\mathbb{E}}
\newcommand{\eps}{\varepsilon}
\newcommand{\dd}{\,\mathrm{d}}
\begin{document}
\section*{UnsolvedMath ID: 20000395}
\textbf{Source problem code:} catalogue entry
\subsection{Definitions and mathematical statement}
John Voight's AIM question asks for the average number of points on the Jacobian of a hyperelliptic curve. The ground field, ensemble, weighting, genus regime, and meaning of average are not specified in the original. In the finite-field interpretation of the report, equation averages, stack-weighted averages, and unweighted curve-isomorphism averages must be distinguished.

Fix an odd prime power \(q\). For \(g\geq1\), let \(\mathcal H_{2g+1,q}\) be the monic squarefree polynomials of degree \(2g+1\) in \(\F_q[T]\), and let \(C_D\) be the smooth projective curve with affine equation \(y^2=D(T)\). Define the uniform equation average
\[
A_{g,q}=\frac1{|\mathcal H_{2g+1,q}|}
\sum_{D\in\mathcal H_{2g+1,q}}\#\operatorname{Jac}(C_D)(\F_q).
\]
The report attributes to Andrade the fixed-\(q\), high-genus asymptotic
\[
A_{g,q}\sim C_q q^g\quad(g\to\infty),\qquad
C_q=\frac1{1-q^{-1}}\prod_P\left(1-\frac1{|P|^2(|P|+1)}\right),
\]
where \(P\) runs over monic irreducible polynomials in \(\F_q[T]\) and \(|P|=q^{\deg P}\). This is a mean in the specified odd-degree equation ensemble, not automatically an average over all hyperelliptic moduli points.

The candidate correction expansion is, as \(q\to\infty\) through odd prime powers,
\[
C_q=1+q^{-1}+q^{-3}+2q^{-6}+O(q^{-7}).
\]
In particular the report asserts that the coefficients of \(q^{-2},q^{-4},q^{-5}\) vanish. This expansion concerns the constants arising after the fixed-field large-genus limit; it is not a uniform joint \((q,g)\) asymptotic unless further estimates are given.

The report also claims a square-term local-density divisor-sum identity for \(C_q\) over monic polynomials \(H\), a stack-groupoid reduction, and an exact genus-one normalization check. The actual summand, stack normalization, and genus-one formula are not supplied in the record, so these cannot serve here as explicit reproducible identities. The report leaves even-degree and characteristic-two ensembles and the unqualified original averaging question distinct.
\subsection{Short English statement}
Specify the averaging ensemble before asking for a Jacobian point-count mean. The report gives an odd-degree equation asymptotic and a large-field expansion of its constant, while retaining separate stack, even-degree, and characteristic-two questions.
\subsection{Sources}
\begin{itemize}
\item AIM workshop ``Arithmetic statistics over finite fields and function fields,'' Question 17, attributed to John Voight: \url{https://aimath.org/pastworkshops/arithstatffieldproblems.pdf}. No PDF was accessed here.
\item Local canonical registry \texttt{aim-algebraic-number-theory-notes.json}, \texttt{notes[46]}, tag \texttt{question}; permanent UnsolvedMath ID \(20000395\), source code \texttt{catalogue entry}. Entire original asks for the average number of points on a hyperelliptic Jacobian. Original remarks and literature field are empty.
\item Supplied research attempt 1 attributes the odd-degree asymptotic to Andrade and proposes the correction expansion and divisor-sum audit. No full authoritative bibliographic references or formulas for the auxiliary checks are supplied. The report is machine generated, not peer reviewed, and requires independent verification before citation.
\end{itemize}
\end{document}
